% Generated draft: requires mathematical and source review.
\documentclass[11pt]{article}
\usepackage[margin=25mm]{geometry}
\usepackage{fontspec}
\setmainfont{Latin Modern Roman}
\usepackage{amsmath,amssymb,mathtools,unicode-math}
\setmathfont{Latin Modern Math}
\usepackage{enumitem,needspace,xurl,hyperref}
\setlength{\parindent}{0pt}
\setlength{\parskip}{5pt}
\setlength{\emergencystretch}{3em}
\newcommand{\N}{\mathbb{N}}
\newcommand{\Z}{\mathbb{Z}}
\newcommand{\Q}{\mathbb{Q}}
\newcommand{\R}{\mathbb{R}}
\newcommand{\C}{\mathbb{C}}
\newcommand{\F}{\mathbb{F}}
\newcommand{\A}{\mathbb{A}}
\newcommand{\PP}{\mathbb{P}}
\newcommand{\E}{\mathbb{E}}
\newcommand{\eps}{\varepsilon}
\newcommand{\dd}{\,\mathrm{d}}
\begin{document}
\section*{UnsolvedMath ID: 86}
\textbf{Source problem code:} catalogue entry
\subsection{Definitions and mathematical statement}
A number field is a finite extension $K/\Q$. For every place $v$ of $K$, write $K_v$ for its completion, and for a finite place write $\mathcal O_v$ for its valuation ring. The idele group $\mathbb A_K^{\times}$ is the restricted product of the groups $K_v^{\times}$ with respect to $\mathcal O_v^{\times}$ at finite places: its elements have a unit component at all but finitely many finite places. Its topology is the restricted-product topology. The idele class group is $C_K=\mathbb A_K^{\times}/K^{\times}$, where $K^{\times}$ is embedded diagonally.

For a finite extension $L/K$, the local norms define an idele norm, whose component at $v$ is the product of $N_{L_w/K_v}(x_w)$ over the places $w$ above $v$. It induces $N_{L/K}:C_L\to C_K$.

Hilbert's ninth problem (1900) asks for a reciprocity law over arbitrary number fields generalizing quadratic reciprocity. Under the supplied literature triage's explicit interpretation in terms of abelian class field theory, a standard precise reciprocity statement is the following. For every number field $K$ and every finite abelian extension $L/K$, there is a continuous surjective Artin reciprocity homomorphism
\[
\theta_{L/K}:C_K\longrightarrow \operatorname{Gal}(L/K),\qquad
\ker\theta_{L/K}=N_{L/K}(C_L),
\]
and hence an isomorphism
\[
C_K/N_{L/K}(C_L)\cong\operatorname{Gal}(L/K).
\]
Normalize this map so that at a finite place $v$ unramified in $L$, the class of an idele with a uniformizer in its $v$ component and $1$ elsewhere maps to the arithmetic Frobenius at $v$, which acts on the residue field by $x\mapsto x^{q_v}$, where $q_v$ is the size of the residue field of $K$ at $v$. Abelianity makes this element independent of the choice of a place above $v$.

This is a precise modern formulation of the abelian reciprocity scope singled out by the supplied assessment, rather than a verbatim historical formulation. ``Arbitrary number fields'' quantifies the base field $K$; it does not assert an analogous norm-quotient theorem for arbitrary nonabelian extensions. The source background mentions Artin's 1927 reciprocity law and continued research on reciprocity. Its dated triage (2026-08-17) classifies the intended historical problem as solved by modern class field theory and explicitly places nonabelian reciprocity programmes outside this statement, whereas the registry status field still says open.
\subsection{Short English statement}
Generalize quadratic reciprocity to every number field. In the supplied abelian class-field-theory interpretation, Artin reciprocity identifies the Galois group of each finite abelian extension with the base field's idele class group modulo norms. The dated source assessment treats this historical problem as resolved and excludes broader nonabelian programmes.
\subsection{Sources}
\begin{itemize}
\item Encyclopedia of Mathematics, ``Hilbert problems,'' accessed 2026-08-17 according to the supplied triage: \url{https://encyclopediaofmath.org/wiki/Hilbert_problems}. The record cites this authoritative secondary source for the classical resolution context.
\item Local registry data/unsolvedmath/problems.json, permanent catalogue entry, source code catalogue entry, ``Hilbert's 9th Problem: Reciprocity Laws,'' attributed to David Hilbert (1900); statement, background, and literature triage dated 2026-08-17. The background attributes the Artin reciprocity law to Emil Artin (1927) but supplies no separate bibliographic citation for it. Dataset attribution: Ulam AI and the UnsolvedMath Contributors.
\end{itemize}
\end{document}
